\documentclass[11pt]{article}

% Shared preamble for the SKJS soundness corpus (main paper + deep dives).
% Loaded via \input in each document's preamble. Documents may define
% additional packages, theorem environments, and macros AFTER this input,
% but must not redefine anything here.
\usepackage{amsmath}
\usepackage{amssymb}
\usepackage{amsthm}
\usepackage{mathtools}
\usepackage{stmaryrd}
\usepackage[margin=1in]{geometry}
\usepackage{tcolorbox}
\usepackage{booktabs}
\usepackage{array}
\usepackage{longtable}
\usepackage[hyphens]{url}
\urlstyle{tt}

% Long typewriter tokens (paths, identifiers) defeat hyphenation;
% allow inter-word stretch instead of overfull lines.
\emergencystretch=2em

\newtheorem{theorem}{Theorem}
\newtheorem{lemma}[theorem]{Lemma}
\newtheorem{definition}[theorem]{Definition}
\newtheorem{corollary}[theorem]{Corollary}
\newtheorem{axiom}[theorem]{Axiom}
\newtheorem{obligation}[theorem]{Obligation}

% Summary box for key definitions
\newtcolorbox{keydef}{
  colback=gray!5,
  colframe=gray!50,
  boxrule=0.5pt,
  arc=2pt,
  left=4pt,right=4pt,top=2pt,bottom=2pt
}

% Core semantic notation shared by every document
\newcommand{\sem}[1]{\llbracket #1 \rrbracket}
\newcommand{\Val}{\mathsf{Val}}

% Implementation anchor: file / class / function at the pinned revision
\newcommand{\impl}[1]{\par\noindent{\footnotesize\textsf{impl:} \texttt{#1}}}

% Single source for the modeled revision. Every document's \date uses this;
% the ledger and tooling read the same values.
\newcommand{\skjspinhash}{9108f339}
\newcommand{\skjspindate}{2026-07-02}
\newcommand{\skjspin}{Models revision \texttt{\skjspinhash} of \texttt{main} (\skjspindate)}



\newtheorem{finding}[theorem]{Finding}
\newtheorem{policy}[theorem]{Policy}

% Macros copied from the main paper (skjs-type-soundness.tex @ \skjspinhash).
\newcommand{\ObjVal}{\mathsf{ObjVal}}
\newcommand{\Env}{\Gamma}
\newcommand{\RefineB}[4]{#1 \vdash #2 \stackrel{#4}{\rightsquigarrow} #3}
\newcommand{\refineop}{\mathsf{refine}}
% Local shorthands.
\newcommand{\Loc}{\mathsf{Loc}}
\newcommand{\decl}{\mathsf{decl}}
\newcommand{\capt}{\mathsf{capt}}
\newcommand{\havoc}{\mathsf{havoc}}

\title{Refinement Invalidation under Mutation\\
{\large Deep Dive F for ``Refinement and Inference Growth in SKJS''}}
\author{}
\date{\skjspin}

\begin{document}
\maketitle

\begin{abstract}
The main soundness paper proves \emph{one-step branch soundness}: a
condition is evaluated against a fixed store, and the refined environment
is guaranteed only at that instant. Its assumption ledger records
\textbf{L-PURE} as \textsc{violated} (issues \#597, \#599) and a tracked
instance of the \textbf{L-MUT} boundary (issue \#598). This deep dive is
part formalization, part design document. It extends the one-step store
model with mutation events, classifies the tracked unsoundness shapes and
pins where each breaks the main paper's branch-soundness argument,
establishes by code inspection and by running the checker what the
implementation does at every mutation site (including two findings beyond
the tracked issues: assignments to captured locals inside closures are
\emph{entirely unchecked}, and closure bodies \emph{read} enclosing
narrowings that may be stale at invocation), analyzes four invalidation
policies, and recommends one — capture-based invalidation for locals plus
field re-widening at writes and call-boundary havoc for field narrowings —
together with the preservation-style lemma the chosen policy validates.
\end{abstract}

\tableofcontents

%==========================================================================
\section{Purpose and Ledger Rows Addressed}
\label{sec:purpose}

This paper addresses the following rows of the main paper's Assumptions
and Proof Obligations table (its \S 8):
\begin{itemize}
\item \textbf{L-PURE} (\textsc{violated}) — ``evaluating a condition does
  not mutate the narrowed locals'', and, by extension, the unstated
  companion assumption that the store does not change between the
  condition and the \emph{uses} of the narrowed environment. Tracked as
  issues \#597 and \#599.
\item \textbf{L-MUT} (\textsc{assumed}, with a tracked instance) — field
  writes are governed by $\mathsf{set}$-types; the tracked instance is
  issue \#598: a field reassignment inside a refined branch does not
  re-widen the narrowed field.
\item \textbf{L-LOCAL} (\textsc{proved} for the module-level half) — the
  captured-local half, which the main paper says ``sharpens L-PURE'', is
  made precise here (Findings~\ref{find:unchecked-write}
  and~\ref{find:stale-read}).
\end{itemize}
The deliverable is a \emph{specified} invalidation policy
(Policy~\ref{pol:main}) with a soundness statement
(Lemma~\ref{lem:preservation}), plus the exact ledger-row edits the main
paper should adopt (Section~\ref{sec:status}). L-PURE remains
\textsc{violated} until the policy lands; this paper turns it from an open
design question into a fix with a stated lemma.

%==========================================================================
\section{Minimal Recap of the Main Paper}
\label{sec:recap}

We reuse the main paper's notation without redefining it: values $\Val$
with runtime tags, object values $\ObjVal$ with the snapshot interface
$(\mathsf{props}, \mathsf{get}, \dots)$, type denotations
$\sem{\cdot}_\rho$ indexed by the substitution state $\rho$, environments
$\Env$ mapping locals to \texttt{TType}, stores $\sigma$ mapping locals to
values, satisfaction $\sigma \models_\rho \Env$, the evaluation relations
$\sigma \models_b e$, and branch soundness for rules
$\RefineB{\Env}{e}{\Env'}{b}$ (main paper, Definitions ``Environments
and Stores'' through ``Branch Soundness''). Three
recap points are load-bearing here:
\begin{enumerate}
\item \textbf{One-step scoping.} The main paper's branch-soundness
  theorems say: if $\sigma \models_\rho \Env$ and $\sigma \models_b e$
  then $\sigma \models_{\rho'} \Env'$ — for the \emph{same} $\sigma$. No
  claim is made about later program points; assignment, loops, and joins
  are covered only by merge soundness (its Lemma ``Merge Soundness'').
  The gap between
  ``$\Env'$ holds at branch entry'' and ``$\Env'$ holds at each use inside
  the branch'' is exactly the subject of this paper.
\item \textbf{Snapshot object semantics.} Membership
  $o \in \sem{\mathsf{object}(B)}$ constrains the properties of $o$ ``at
  the moment a condition is evaluated'' (main paper, Definition
  ``Object Interface''), via
  $\mathsf{get}$-types only; $\mathsf{set}$-types govern which writes the
  checker accepts (\textbf{L-MUT}) and do not participate in membership.
\item \textbf{The refinement primitive writes to the scope.} All narrowing
  lands in the flow-sensitive scope by
  \texttt{setLocalFilterMap}/\texttt{scope.set}; field narrowing
  (\textsc{R-EqField}, \textsc{R-TruthyField}) rewrites the field's
  $\mathsf{get}$ inside the local's object type via
  \texttt{narrowField} $\to$ \texttt{updateField} $\to$
  \texttt{TField::withGet}, leaving $\mathsf{set}$ untouched.
\end{enumerate}
\impl{TypeCheckRefine.sk / setLocalFilterMap, narrowField;
TType.sk / TField::withGet}

\paragraph{The scope chain.}
Because the invalidation question is about \emph{where} flow information
lives, we fix the implementation's scope discipline once. A
\texttt{Scope} is a mutable frame with an \texttt{up} pointer, a
per-frame \texttt{isLoop} flag, declared names (\texttt{locals}), and
flow entries (\texttt{refs}). Lookup (\texttt{maybeGetRef}) walks the
\texttt{up} chain and records whether it crossed a loop frame. The
central asymmetry: \texttt{Scope::set} \emph{always writes into the
current frame}, even when the binding was found higher up — flow updates
shadow rather than mutate the defining frame — and if the lookup crossed
a loop frame, the assigned type is first checked against the type the
binding had at loop entry (the back-edge guard). Structured constructs
reconcile shadow frames explicitly: \texttt{if}/\texttt{else} via
\texttt{mergeScope} (pointwise union) or \texttt{popScopeReplace} when
one side terminates; loop and \texttt{try} frames via
\texttt{popScope} (union with the parent entry). Closure bodies are the
exception and the subject of Section~\ref{sec:sites}.
\impl{TypeCheckScope.sk / Scope::set, maybeGetRefHelper, merge, popMerge,
popReplace; TypeCheckStmt.sk / If, While, For arms}

%==========================================================================
\section{A Store Model with Mutation Events}
\label{sec:model}

The main paper's store $\sigma$ maps locals to values and is fixed. To
talk about invalidation we split it and let it step.

\begin{definition}[Two-Level Store]
\label{def:store}
A store is a pair $\sigma = (\sigma_L, \sigma_H)$ where
$\sigma_L : \mathsf{Var} \rightharpoonup \Val$ maps locals to values and
$\sigma_H$ is a heap giving each object identity its current property
map; the snapshot interface $\mathsf{get}(o, p)$ of the main paper's
the Object Interface definition is read as
$\mathsf{get}_{\sigma_H}(o, p)$. Satisfaction
$\sigma \models_\rho \Env$ means
$\sigma_L(x) \in \sem{\Env(x)}^{\sigma_H}_\rho$ for all
$x \in \mathrm{dom}(\Env)$, where object-type membership consults
$\sigma_H$. All results of the main paper hold verbatim for a fixed
$(\sigma_L, \sigma_H)$.
\end{definition}

\begin{definition}[Mutation Events]
\label{def:events}
Between two checker-visible program points the store may step
$\sigma \xrightarrow{\mu} \sigma'$ by one of:
\begin{itemize}
\item \textbf{(M-Asgn)} direct assignment $x := v$: replaces
  $\sigma_L(x)$; $\sigma_H$ unchanged.
\item \textbf{(M-Field)} field write $o.f := v$: replaces one property in
  $\sigma_H$; $\sigma_L$ unchanged. Reachable from the checked function as
  \texttt{x.f = e}, or from \emph{any} called code holding an alias of
  $o$. If $f$ is accessor-backed, the event is instead ``run the setter'',
  which may perform arbitrary further events.
\item \textbf{(M-Call)} call/new/await boundary: runs foreign code, which
  may perform any sequence of (M-Field) events on heap objects reachable
  from its arguments, from globals, or from closures it can reach, and any
  (M-Asgn) events on locals \emph{captured by such closures}. Locals not
  captured by any closure are unreachable: (M-Call) cannot write them.
\end{itemize}
JS is single-threaded per agent, so between two adjacent checker-visible
points with no intervening call and no intervening syntactic write, the
store is genuinely fixed; this is what makes the main paper's one-step
theorems non-vacuous.\footnote{In a multi-agent (shared-memory) setting
even that fails for \texttt{SharedArrayBuffer}-backed state; SKJS, like
tsc, ignores this, and so do we.}
\end{definition}

\begin{definition}[Preservation Discipline]
\label{def:discipline}
The checker tracks a current environment $\Env_t$ at each program point
$t$. Call an environment-update function $U_\mu$ (what the checker does
to $\Env$ when it sees, or must conservatively assume, event $\mu$)
\emph{adequate} for $\mu$ if
$\sigma \models_\rho \Env$ and $\sigma \xrightarrow{\mu} \sigma'$ imply
$\sigma' \models_{\rho}U_\mu(\Env)$. Flow-sensitive checking is sound iff
every event that can occur between two points is matched by an adequate
update on that path. The main paper's branch-soundness results supply the
$U$ for condition evaluation (no event, $\sigma' = \sigma$); this paper
supplies — or exhibits the absence of — the $U$'s for
(M-Asgn), (M-Field), (M-Call).
\end{definition}

\subsection{The three tracked shapes}
\label{sec:shapes}

Each shape below is pinned by
\texttt{tests/expect-error/unsound-refinement.ts} (cases 1--3); all three
compile without error today, re-confirmed by running the checker at the
pinned revision.

\paragraph{Shape (a) — \#597: side-effecting call between test and use.}
\begin{verbatim}
declare function sideEffect(): void;
function f(x: { f: number | null }): void {
  if (x.f !== null) {
    sideEffect();          // (M-Call): may write x.f through an alias
    const n: number = x.f; // accepted; unsound
  }
}
\end{verbatim}
\textsc{R-EqField} narrows the object branch's field to
$f.\mathsf{get} = \texttt{number}$ at branch entry — sound there. The
call is an (M-Call) event: $x$'s object is reachable from the caller of
$f$ (it is a parameter, hence aliased), so $\sigma_H$ may step. The
checker's update for the call is the identity, which is \emph{not}
adequate for (M-Call): the failing hypothesis is the fixity of $\sigma_H$
— precisely assumption \textbf{L-PURE} read as ``the snapshot taken by
the snapshot Object Interface is still current at the use''. No lemma of
the main paper is
wrong; its scope boundary (``one-step'') is real.

\paragraph{Shape (b) — \#599: captured local mutated via an invoked lambda.}
\begin{verbatim}
function f(): void {
  let x: string | null = "hello";
  if (x !== null) {
    const reset = () => { x = null; };
    reset();               // (M-Call) performing (M-Asgn) on captured x
    const s: string = x;   // accepted; unsound
  }
}
\end{verbatim}
Here the mutation is not foreign: the checker \emph{typechecks the
assignment} \texttt{x = null} while checking the lambda body. The event
is visible, yet the update applied to the branch environment is the
identity — Section~\ref{sec:sites} shows the assignment lands in the
lambda's scope frame, which is discarded. The failing hypothesis is again
\textbf{L-PURE} at the \texttt{reset()} boundary; but unlike shape (a),
an exact (not conservative) update was available and thrown away.

\paragraph{Shape (c) — \#598: field reassignment inside a refined branch.}
\begin{verbatim}
function f(x: { f: number | null }): void {
  if (x.f !== null) {
    x.f = null;            // (M-Field), checked against set-type: passes
    const n: number = x.f; // accepted; unsound (tsc rejects, TS2322)
  }
}
\end{verbatim}
A pure (M-Field) event with no call and no alias needed: the write is
syntactic and checked — against the field's $\mathsf{set}$-type
(\texttt{number | null}, untouched by narrowing), so it passes — but the
update applied to $\Env(x)$ is the identity, leaving the stale
$f.\mathsf{get} = \texttt{number}$. The failing hypothesis is the
\textbf{L-MUT} note's ``membership semantics (get-only) is sound for the
snapshot reads refinement performs'': after the write the snapshot is
stale, so $\sigma'(x) \notin \sem{\Env(x)}^{\sigma'_H}$. This shape is
\emph{not} tsc-parity: tsc's assignment narrowing rewrites the declared
flow type of \texttt{x.f} to \texttt{null} and rejects the next line.
\impl{TypeCheckRefine.sk / evalFieldEqualityType (withGet);
TType.sk / TFields::resolveFieldType (isLValue reads set)}

%==========================================================================
\section{What the Implementation Does at Each Mutation Site}
\label{sec:sites}

Established by code inspection at \texttt{\skjspinhash} and, where marked
\emph{(probed)}, by running the checker (\texttt{make
check-file-in-docker}) on minimal programs.

\subsection{(M-Asgn) direct assignment to a local: sound flow update}
\texttt{assign}'s \texttt{Var}-LHS arm types the RHS and calls
\texttt{scope.set(n, ty2, ...)}: the local's flow entry is
\emph{replaced} by the RHS type in the current frame. This is the
adequate update $U_{\text{M-Asgn}}(\Env) = \Env[x \mapsto \tau_{\rm RHS}]$
($\sigma'_L(x) = v \in \sem{\tau_{\rm RHS}}$ by \textbf{L-TC}). In
particular a narrowed type is correctly clobbered: \texttt{if (x !== null)
\{ x = null; const s: string = x; \}} is rejected \emph{(probed)}. Through
a loop frame, \texttt{set} additionally checks
$\tau_{\rm RHS} <: {}$loop-entry type — the back-edge guard that keeps
iteration-2 reads justified; narrowed loop-entry types participate:
narrowing $x$ to \texttt{string} before a \texttt{while}, then assigning
\texttt{null} inside it, is rejected \emph{(probed)}.
\impl{TypeCheckExpr.sk / assign (Var arm, $\sim$line 3016);
TypeCheckScope.sk / Scope::set (throughLoop check)}

\subsection{(M-Asgn) through a closure: unchecked and invisible}
\texttt{anonFunction} checks a lambda body in a fresh subscope whose
\texttt{up} is the creation-site scope, then restores the saved scope
\emph{without any merge}: the lambda frame and every flow entry written
into it are discarded. Combined with the \texttt{Scope::set} asymmetry
(writes go to the current frame; the loop-guard subtype check fires only
across \texttt{isLoop} frames, and lambda frames are not loop frames),
assignments to captured locals are (i) never propagated and (ii) never
checked against the outer binding's type. Two consequences beyond the
tracked issues:

\begin{finding}[Captured writes are entirely unchecked]
\label{find:unchecked-write}
\emph{(probed)} The program
\texttt{let x: string = "a"; const g = () => \{ x = 5; \}; g();
console.log(x.length);}
compiles without any diagnostic; at runtime \texttt{x.length} is
\texttt{undefined}. This is stronger than \#599: no refinement is
involved, and the type confusion is arbitrary (\texttt{number} into
\texttt{string}), not merely a stale narrowing. tsc rejects the
assignment (TS2322). Verified distinct from every tracked issue
(\#597--\#599, \#585) and confirmed with a runtime witness; issue
draft prepared:
\path{issues/captured-local-writes-unchecked.md}. Ledger row proposed
in Section~\ref{sec:status}.
\end{finding}

\begin{finding}[Closure bodies read narrowings that are stale at invocation]
\label{find:stale-read}
\emph{(probed)} Because lookup walks the whole \texttt{up} chain
(\textbf{L-LOCAL}'s captured-local half), a lambda body checked inside a
refined branch \emph{sees the narrowed type}:
\texttt{let x: string | null = "hi"; if (x !== null) \{ const f = () =>
x.length; x = null; f(); \}}
compiles clean and crashes at runtime. The narrowing was valid when the
closure was \emph{created}, not when it \emph{runs}. tsc handles this by
re-widening, inside closure bodies, any \texttt{let} that is assigned
anywhere in the function. This is the dual of shape (b): (b) is a stale
write-out, this is a stale read-in. Confirmed with a runtime witness (a \texttt{TypeError} reading
\texttt{.length} of \texttt{null}); verified not covered by \#599's
write-out text; issue draft prepared:
\path{issues/stale-closure-reads-narrowing.md}.
\end{finding}
\impl{TypeCheckExpr.sk / anonFunction ($\sim$line 2850: scope
save/restore, no popScope); TypeCheckScope.sk / maybeGetRefHelper}

\subsection{(M-Field) syntactic field write: checked, but no update}
For \texttt{x.f = e} (with \texttt{x} not a tracked growable local
object), \texttt{assign} falls to its default arm: it types the LHS in
lvalue mode — \texttt{dot} with \texttt{AKLValue}, which resolves the
field via \texttt{resolveFieldType(.., isLValue=true)} $=$ the field's
$\mathsf{set}$-type, intersected across members — checks
$\tau_{\rm RHS} <:$ that, and returns. \emph{It never touches the scope
entry for \texttt{x}}: there is no \texttt{scope.set} on this path, so a
narrowed $f.\mathsf{get}$ survives its own overwrite (shape (c)). The
narrowing machinery and the write path even share the data: narrowing
rewrote $f.\mathsf{get}$ via \texttt{withGet} and left $f.\mathsf{set}$
declared, which is exactly why the write still typechecks. (The separate
\texttt{assignLocalObject} path for growable local objects \emph{does}
update the scope, by intersection — a growth mechanism, not an
invalidation one, and out of scope here.)
\impl{TypeCheckExpr.sk / assign (default arm), dot (AKLValue);
TType.sk / TFields::resolveFieldType}

\subsection{(M-Call): no invalidation whatsoever}
No call, \texttt{new}, \texttt{await}, or method invocation touches
\texttt{tcenv.scope}. The complete set of scope writers in the checker is:
\texttt{scope.set} in the refinement primitives, in \texttt{assign}'s
\texttt{Var} arm, in \texttt{assignLocalObject}, and in the
\texttt{this}-retyping after \texttt{super(...)}; plus the structured
merges (\texttt{mergeScope}/\texttt{popScope}/\texttt{popScopeReplace}).
Call typing is absent from that list, so shapes (a) and (b) follow
mechanically: the identity update is applied at every (M-Call).
\impl{TypeCheckRefine.sk; TypeCheckExpr.sk / assign;
TypeCheckStmt.sk (grep: no scope writes in call paths)}

%==========================================================================
\section{Design Space}
\label{sec:design}

Notation: for a function under check, let $C$ be the set of locals
captured by at least one closure, $W \subseteq C$ those additionally
assigned \emph{anywhere in the function} (matching tsc's discipline and
the stale-read witness, whose write is in the enclosing body, not a
closure; both computable syntactically in one pre-pass), and say a local
is
\emph{heap-visible} if its value is an object that may be aliased —
conservatively: any object-typed local other than a fresh, never-escaping
literal. $\havoc_S(\Env)$ re-widens the entries of $S$: whole-variable
narrowings back to the binding's declared type, and field-level
narrowings by restoring each narrowed $f.\mathsf{get}$ (mechanisms in
Section~\ref{sec:rewiden}).

\subsection{Option (i): havoc-on-call}
At every (M-Call) boundary apply
$\havoc_{W \,\cup\, \{\text{field narrowings of heap-visible locals}\}}$.
\emph{Soundness:} adequate for (M-Call) by construction — the havocked
entries over-approximate every value the event can install (captured
writes are bounded by the declared type \emph{once
Finding~\ref{find:unchecked-write} is fixed}; field writes by
$\mathsf{set}$-types via \textbf{L-MUT}), and unhavocked entries are
unreachable by the event (Definition~\ref{def:events}: uncaptured locals
cannot be written by a call). \emph{Precision cost:} the breaking shape
is ``narrow once, then several statements each containing a call and a
later use of the narrowing''. Spot-checks against the repository corpus
at the pinned revision:
\begin{itemize}
\item \texttt{examples/src/type\_refinement\_examples.ts}
  (\texttt{processResponse}): \texttt{response.kind === "success"}
  narrowing used once, \emph{inside the argument} of the following
  \texttt{console.log} — argument evaluation precedes the call, so a
  call-boundary havoc does not break it. Survives.
\item \texttt{examples/src/deck.ts} (display loop): \texttt{card}
  narrowed by \texttt{!== undefined}, used only in the argument of one
  \texttt{console.log}. Survives.
\item \texttt{examples/src/runtime\_demo.ts} /
  \texttt{builtins\_documentation.ts}: \texttt{aliceAge !== undefined}
  then \texttt{aliceAge.toString()} inside one call argument. Survives —
  but note \texttt{aliceAge} is not captured and not object-typed, so
  option (i) as scoped above would not havoc it anyway; only a cruder
  ``havoc everything'' would, and even that survives these single-use
  sites.
\item Breaking witness (crude variant): any
  \texttt{if (x) \{ log(x.a); log(x.b); \}} with \texttt{x} a narrowed
  \emph{field path or captured local} — the first \texttt{log} havocs,
  the second use fails. \texttt{examples/src/try\_catch.ts} (uses
  \texttt{e.message} after an earlier statement in the branch) is the
  nearest real instance; it survives the scoped option (i) only because
  the intervening statements contain no call.
\end{itemize}
The corpus tolerates the \emph{scoped} havoc well: narrowed values are
overwhelmingly consumed before or inside the first following call. The
cost concentrates on ADT-style field narrowings used across several
statements with interleaved calls — common in real code, rarer in this
repo.

\subsection{Option (ii): effect / purity tracking}
Havoc only at calls whose callee may write the relevant location. SKJS
already builds a call graph during checking (\texttt{CGNode},
\texttt{recordCalls}), so the skeleton exists; what is missing is
(1) a per-function summary ``writes captured local $x$'' /
``writes some field $f$ through some alias'', computed to a fixpoint over
the call graph including indirect calls through function-typed values —
which in practice degrades to ``any unknown callee may write anything'',
(2) summaries for \texttt{declare}d ambient functions, which have no
bodies — so either an effect annotation surface is added to the language
(a major design commitment tsc has repeatedly declined) or ambient calls
havoc everything, and (3) soundness of the summary under first-class
functions, requiring the call graph itself to be sound, which it is not
required to be today (it exists for diagnostics/codegen ordering).
\emph{Verdict:} strictly dominates (i) in precision when it fires, but
its sound version collapses to (i) at every ambient or indirect call —
the overwhelmingly common case (\texttt{console.log} is ambient). Not
recommendable as the primary mechanism; viable later as a precision
overlay on top of (i)/(iii).
\impl{Reachability.sk / CGNode; TypeCheckEnv.sk / recordCalls}

\subsection{Option (iii): capture-based invalidation}
Precompute $W$ (captured-and-assigned locals). Then: (1) at every
(M-Call) boundary, $\havoc_W$; (2) inside closure bodies, do not let
narrowings of $W$-locals flow in — resolve them at their declared type
(fixing Finding~\ref{find:stale-read}); (3) check captured writes against
the declared type (fixing Finding~\ref{find:unchecked-write}, and making
``declared type'' a sound havoc target). This is sound for
\emph{whole-local} narrowings: an (M-Call) can write exactly the
$W$-locals (Definition~\ref{def:events}), and each write is now bounded
by the declared type. It deliberately does \emph{not} address field
narrowings (shape (a)) — the heap is aliased regardless of capture. tsc
implements a cousin of this policy in restricted positions (narrowings of
\texttt{var}s and of \texttt{let}s assigned in nested functions are
dropped inside closure bodies, and loop-carried narrowings are reset),
but performs no call-boundary invalidation at all.
\emph{Precision cost:} negligible in the corpus — no repository example
narrows a captured-and-assigned local and then calls; the pinned
\#599 reproducer is the canonical instance and is exactly what should
break.

\subsection{Option (iv): status quo (documented tsc-parity unsoundness)}
Keep everything, document L-PURE as violated (done), pin the
reproducers (done). This is genuine tsc parity for shapes (a) and (b) —
tsc accepts both — but \emph{not} for shape (c) or
Finding~\ref{find:unchecked-write}, both of which tsc rejects. Status quo
is therefore not a coherent resting point: two of the five behaviors are
strictly-more-unsound-than-tsc bugs, not policy choices.

\subsection{Recommendation}
\label{sec:recommend}

\begin{policy}[Adopted invalidation policy]
\label{pol:main}
Layered, in decreasing order of urgency:
\begin{enumerate}
\item \textbf{P1 (bug fix):} check captured writes against the captured
  binding's declared type (Finding~\ref{find:unchecked-write});
  concretely, \texttt{Scope::set} applies its subtype check when the
  lookup crossed a \emph{function} frame, exactly as it already does for
  loop frames.
\item \textbf{P2 (\#598):} re-widen the narrowed field getter at every
  syntactic field write (rule in Section~\ref{sec:rewiden}).
\item \textbf{P3 (\#599 + Finding~\ref{find:stale-read}):} option (iii)
  capture-based invalidation for $W$-locals: havoc at calls, block
  narrowing inflow to closure bodies.
\item \textbf{P4 (\#597):} option (i) scoped to field-level narrowings of
  heap-visible locals, applied at (M-Call) boundaries. This is the only
  layer with a real precision cost; if product priorities delay it, the
  deviation is exactly tsc's, and L-PURE's ledger row shrinks to this one
  named residue rather than three shapes.
\end{enumerate}
\end{policy}

\begin{lemma}[Preservation across mutation events — replaces one-step
scoping at mutation sites]
\label{lem:preservation}
Assume the main paper's axioms (\textbf{L-UNION}, \textbf{L-INTER},
\textbf{L-EXP}, \textbf{L-MONO}, \textbf{L-TC}, \textbf{L-SUB}) and the
well-formedness condition $\sem{f.\mathsf{get}} \supseteq$ the values a
read of $f$ can yield after an accepted write
(Obligation~\ref{obl:getset}). Under Policy~\ref{pol:main} (P1--P4), for
every checker path point $t$ with environment $\Env_t$ and every event
$\sigma \xrightarrow{\mu} \sigma'$ occurring between $t$ and the next
point $t'$: $\sigma \models_\rho \Env_t$ implies
$\sigma' \models_\rho \Env_{t'}$, where $\Env_{t'}$ is the environment
produced by the policy's update for $\mu$
(Definition~\ref{def:discipline}):
$U_{\text{M-Asgn}} = [x \mapsto \tau_{\rm RHS}]$ with the P1 declared-type
check for captured writes; $U_{\text{M-Field}}$ the P2 re-widening;
$U_{\text{M-Call}} = \havoc_W \circ \havoc_{\rm fields}$ (P3, P4).
Consequently branch soundness extends from branch entry to every use in
the branch, by induction over the events on the path.
\end{lemma}

\noindent\emph{Proof sketch.} Case (M-Asgn): \textbf{L-TC} for the RHS;
for captured writes the P1 check bounds the written value by the declared
type, which is what $\havoc$ later restores. Case (M-Field):
Lemma~\ref{lem:rewiden-sound}. Case (M-Call): the event can write only
$W$-locals and the heap (Definition~\ref{def:events}); $\havoc_W$
restores declared types, sound by the P1 bound; $\havoc_{\rm fields}$
restores get-types whose adequacy is Obligation~\ref{obl:getset};
entries outside both havoc sets denote locations the event cannot write,
so their satisfaction transports unchanged. $\square$

%==========================================================================
\section{The Re-Widening Rule for Field Writes (\#598)}
\label{sec:rewiden}

\begin{keydef}
\textbf{Rule W-Field.} On \texttt{x.f = e} with RHS type $\tau_{\rm RHS}$
(after the existing $\tau_{\rm RHS} <: \bigsqcap f.\mathsf{set}$ check),
update \texttt{x}'s scope entry by the memberwise primitive
$\refineop^{g,\mathsf{None}}$ with mapper: non-object members unchanged;
each object branch $\beta$ with a declared field $f$ becomes
\[
\beta[f.\mathsf{get} \mapsto
  \mathsf{union}(f.\mathsf{get},\ \tau_{\rm RHS}),\quad
  f.\mathsf{opt} \mapsto f.\mathsf{opt}]
\]
and branches without $f$ (or reached only via an index) are unchanged.
Optionality is \emph{kept}: clearing it would be sound for data fields
(the write makes $f$ present) but unsound for a branch whose write went
through a setter-less index — keep the conservative reading.
\end{keydef}

This reuses the exact machinery of \texttt{narrowField}: the same
\texttt{setLocalFilterMap} walk, the same
\texttt{updateField}$/$\texttt{withGet} write-back — with
$\mathsf{union}$ where narrowing used $\mathsf{inter}$/$\mathsf{rmc}$,
and with the branch-dropping arm removed (a write never refutes a
branch). No new representation is needed.
\impl{TypeCheckRefine.sk / narrowField, setLocalFilterMap (mechanism to
reuse); TypeCheckExpr.sk / assign (default arm — the missing
scope.set site)}

\begin{lemma}[W-Field is adequate for (M-Field)]
\label{lem:rewiden-sound}
If $\sigma \models_\rho \Env$ and $\sigma \xrightarrow{o.f := v} \sigma'$
with $v \in \sem{\tau_{\rm RHS}}$ and $o = \sigma_L(x)$, then
$\sigma' \models_\rho U(\Env)$ for the W-Field update $U$ — provided $f$
is data-backed in $o$ \emph{and} no other local's recorded type views the
same object through a wider $\mathsf{set}$-type for $f$ (cross-alias
get/set coherence — registry row \textbf{L-MUT-ALIAS}). The latter
hypothesis fails at the pin: the alias
\texttt{const y:\ \{f:\ number\} = x} \emph{is admitted} after narrowing
and the legal write \texttt{o.f := null} through \texttt{x} is
unrecorded on $y$ (mechanism, executed witness, and tsc comparison:
issue draft \path{issues/narrowed-alias-stale-write.md}). The
hypothesis is strictly stronger than Obligation~\ref{obl:getset}; its
site-by-site grading — including the minted views that need no declared
alias at all — is deep dive H's.
\end{lemma}
\begin{proof}
Only $x$'s entry changes and only $o$'s $f$-property changed in
$\sigma_H$. Let $\beta$ be the object branch with $o \in \sem{\beta}$.
If $\beta$ declares $f$: every constraint of $\beta$ except the $f$-value
constraint is insensitive to the event; the new value $v$ satisfies the
new get-type by $v \in \sem{\tau_{\rm RHS}} \subseteq
\sem{\mathsf{union}(f.\mathsf{get}, \tau_{\rm RHS})}$ (\textbf{L-UNION}).
If $\beta$ does not declare $f$ but an index covers it, the index
constraint on the new value holds because the accepted write checked
$\tau_{\rm RHS}$ against the resolved $\mathsf{set}$-type, which for an
index-covered key includes the index's set-type, and the well-formedness
of Obligation~\ref{obl:getset} carries it to the get side. Entries of
other locals may alias $o$: their branches' $f.\mathsf{get}$ must also
admit $v$ — this is \emph{not} guaranteed by W-Field (it updates only
$x$) and is exactly why P4 additionally havocs field narrowings of other
heap-visible locals at writes-through-aliases; for the syntactic
single-local write of \#598, aliasing within the same function without an
intervening call also requires havocking \emph{other} locals' narrowed
$f$ — the implementation should apply W-Field to every local whose type
carries a narrowed field named $f$, a cheap over-approximation.
\end{proof}

\paragraph{The accessor situation.}
At the pinned revision SKJS does not represent \texttt{get}/\texttt{set}
accessor members as get/set fields at all: the parser records an
\texttt{isGetSet} marker (\texttt{FGet}/\texttt{FSet}) but the type
checker never consults it, so a getter \texttt{get p(): T} is typed as a
\emph{method} \texttt{p(): T} (a setter likewise, then rejected for
lacking a return type). That method-typing is itself unsound in the
\emph{call} direction --- \texttt{const n: number = c.p()} type-checks
but throws at runtime, and tsc rejects it (TS6234); registry row
ACCESSOR-METHOD (issue draft \path{getter-typed-as-method.md}), a known
limitation (tests \path{callgraph-unsupported/getter-calls.ts}). For
re-widening the consequence is that there are no coercing-setter fields:
a getter read is never a narrowable value (it is a function), and the
only fields with genuinely distinct \texttt{get}/\texttt{set} are
synthetic index signatures (independent get/set variables) and readonly
fields (\texttt{set}${=}\bot$), neither of which runs user code on a
write. So the W-Field rule's $\mathsf{union}$ against $\tau_{\rm RHS}$
carries no accessor residue at present; its only residue is the
cross-alias one of Lemma~\ref{lem:rewiden-sound}.

\paragraph{The accessor caveat, forward-looking.}
\emph{If} accessor members were represented as get/set fields, writing
\texttt{o.f = e} would run a coercing setter, which could store --- and a
subsequent get return --- any value of the \emph{declared} get-type,
which $\mathsf{union}(\text{narrowed get}, \tau_{\rm RHS})$ need not
contain. The sound re-widening would then be
$\mathsf{union}(f.\mathsf{get}, f.\mathsf{set})$, which contains the
declared get-type whenever declared-get $\subseteq$ set holds
(Obligation~\ref{obl:getset}); tsc's $\tau_{\rm RHS}$ narrowing is unsound
for coercing accessors in the same way. Because \texttt{TField} records
only \texttt{get}/\texttt{set}/\texttt{access} and not accessor-ness, the
$\mathsf{union}(f.\mathsf{get}, f.\mathsf{set})$ form would be the sound
option expressible even then --- and supplying accessor-ness, so the
checker builds get/set accessor fields, is the same change that closes
ACCESSOR-METHOD. We therefore recommend the $\tau_{\rm RHS}$ variant now
(sound today up to the cross-alias residue, and matching tsc), revisiting
the accessor case if and when accessor fields are added.

\begin{obligation}[Get/set coherence]
\label{obl:getset}
For every well-formed field, any value readable from the field after an
accepted write lies in $\sem{f.\mathsf{get}}$ read at the field's
\emph{declared} (un-narrowed) type. For data fields this is immediate
(the read returns the written value, bounded by
$f.\mathsf{set} = f.\mathsf{get}$ declared); for accessors it needs the
declared-get $\subseteq$ set check, which SKJS should enforce at
declaration (tsc's TS2380 lineage). Status: \textsc{open}; needed by
Lemma~\ref{lem:preservation}'s field case and by exit (1) above.
\end{obligation}

%==========================================================================
\section{Status and Registry Impact}
\label{sec:status}

The authoritative, current statuses are maintained in the registry
(\texttt{docs/ledger.yaml}), rendered as the main paper's Assumptions and
Proof Obligations tables; this paper is cited there as evidence. The
registry rows it supports:

% GENERATED by docs/tools/docs.py from docs/ledger.yaml — do not edit.

\begin{center}
\renewcommand{\arraystretch}{1.25}
\small
\begin{tabular}{@{}>{\raggedright\arraybackslash}p{2.4cm}>{\raggedright\arraybackslash}p{8.2cm}>{\raggedright\arraybackslash}p{2.6cm}@{}}
\toprule
\textbf{Row} & \textbf{Statement} & \textbf{Status} \\
\midrule
L-LOCAL & Module-level bindings are never narrowed; function-level locals are, including captured ones. & \textsc{proved} \\
L-PURE & Evaluating a condition does not mutate the narrowed locals. & \textsc{violated} \\
L-CAPT & Assignments to captured annotated or initializer-inferred locals are checked against the binding's declared or widened inferred type. & \textsc{proved} \\
L-CAPT-EVOLVE & A captured write to an evolving binding (bare \texttt{let x = null/undefined}) is unchecked, and later reads keep the stale evolved type. & \textsc{violated} \\
L-STALE-READ & A closure body invoked later reads enclosing narrowings that are stale at invocation time. & \textsc{violated} \\
L-LOOP-CARRIED & A pre-loop narrowing of an annotated binding survives a back-edge write that invalidates it. & \textsc{violated} \\
L-GETSET & Get/set coherence of accessor fields. & \textsc{open} \\
ACCESSOR-METHOD & Accessor (\texttt{get}/\texttt{set}) members must be typed as fields, not methods; today a getter is a call signature, so \texttt{c.p()} type-checks and throws. & \textsc{violated} \\
\bottomrule
\end{tabular}
\end{center}


\paragraph{Imports (review contract).}
The named statements this paper assumes from elsewhere. A reviewer treats
each as an axiom, checking citations against the registry's canonical
statement; no other paper's internals are in scope:

% GENERATED by docs/tools/docs.py from docs/ledger.yaml — do not edit.

\begin{center}
\renewcommand{\arraystretch}{1.25}
\small
\begin{tabular}{@{}>{\raggedright\arraybackslash}p{2.6cm}>{\raggedright\arraybackslash}p{7.4cm}>{\raggedright\arraybackslash}p{1.2cm}>{\raggedright\arraybackslash}p{2.2cm}@{}}
\toprule
\textbf{Row} & \textbf{Canonical statement} & \textbf{Owner} & \textbf{Status} \\
\midrule
L-UNION & \texttt{union} over-approximates set union (all levels). & B & \textsc{proved} \\
L-INTER & \texttt{inter} over-approximates set intersection (the containment clause). & B & \textsc{proved} \\
L-EXP & Expansion preserves or enlarges semantics (modes \texttt{None}/Covariant); residual instances keep theirs. & C & \textsc{proved} \\
L-MONO & Refinement-time state evolutions preserve membership in already-installed types. & E & \textsc{proved} \\
L-SUB & \texttt{Clauses::isSub} acceptance implies semantic inclusion under the current state. & D & \textsc{open} \\
L-TC & Checked expressions inhabit their checked types. & main & \textsc{assumed} \\
L-MUT & Field writes are governed by $\mathsf{set}$-types; get-only membership is sound for snapshot reads. & main & \textsc{assumed} \\
\bottomrule
\end{tabular}
\end{center}


\noindent What this paper establishes per row, and how the specified
fixes change the registry when they land:

\begingroup
\renewcommand{\arraystretch}{1.3}
\small
\begin{longtable}{@{}p{2.0cm}p{8.6cm}p{3.4cm}@{}}
\toprule
\textbf{Row} & \textbf{Established here / fix impact} &
\textbf{Trajectory when the fix lands} \\
\midrule
\endhead
L-PURE & Fix specified: Policy P1--P4. Lemma~\ref{lem:preservation}
  discharges the \emph{caller-path} shapes (a) \#597 by P4 and (b)
  \#599 by P3. Three residues it does \emph{not} yet cover: the closure
  \emph{stale-read} (Finding~\ref{find:stale-read}) is closed by P3's
  static inflow rule, not by the preservation induction, and needs a
  separate closure-body lemma; a narrowed \emph{field} read inside a
  closure is closed by no layer (a witnessed shape, issue draft
  \texttt{stale-closure-reads-narrowing.md}); and implicit invocations
  (getters/coercions/iterators) are outside the (M-Call) event model, so
  P3/P4 do not fire there. &
  the caller-path shapes close once P3/P4 land; the
  closure-body read and event-model residues remain \\
L-MUT & \#598 has a specified re-widening rule (W-Field, adequate by
  Lemma~\ref{lem:rewiden-sound} for data fields); accessor residue
  gated on get/set coherence (Obligation~\ref{obl:getset}); the
  cross-view facet is the separate row L-MUT-ALIAS, witnessed via
  \texttt{instanceof Array} and via narrowing-admitted aliases, both
  outside W-Field's reach (issue drafts
  \texttt{instanceof-array-mutable-top-view.md},
  \texttt{narrowed-alias-stale-write.md}; graded in deep dive H). &
  data fields discharged once W-Field lands; L-MUT-ALIAS needs its
  own fix (\texttt{ReadonlyArray}) \\
L-LOCAL & The captured-local half made precise: closure bodies read
  enclosing narrowings valid at creation, not invocation
  (Finding~\ref{find:stale-read}); P3 blocks inflow for
  captured-and-assigned locals. &
  captured half governed by L-PURE/L-CAPT \\
L-CAPT & Made precise, with a witness
  (Finding~\ref{find:unchecked-write}): the captured write is unchecked
  and discarded; fix is P1; issue draft
  \path{captured-local-writes-unchecked.md}. &
  closes once P1 lands \\
L-GETSET & Introduced here (Obligation~\ref{obl:getset}):
  prerequisite for the sound accessor variant of W-Field. &
  gated on the get/set coherence check \\
\bottomrule
\end{longtable}
\endgroup

\paragraph{Recommended pinned tests.}
Extend \texttt{tests/expect-error/unsound-refinement.ts} (cases 1--3
already pin shapes (a)--(c)) with:
case 4, Finding~\ref{find:unchecked-write} (captured write of a
wrong-typed value, no refinement involved); case 5,
Finding~\ref{find:stale-read} (closure reads a narrowing, local
reassigned, closure invoked). Each policy layer that lands flips its
cases from expect-clean to expect-error, forcing the corresponding
passage of this paper and of the main paper's ledger to be updated in the
same change — the maintenance discipline of the main paper's
Appendix B. Additionally pin the two \emph{sound} behaviors this paper
leans on, so regressions surface: direct-assignment clobbering
(\texttt{x = null} after narrowing rejects the next use) and the
through-loop back-edge guard (assignment inside a loop checked against
the narrowed loop-entry type) — both verified by probe at the pinned
revision but not currently pinned by any test.

\paragraph{Failed attempts and negative findings, for the record.}
(1) A pure ``havoc everything at every call'' baseline was considered and
rejected without implementation: spot-checks (Section~\ref{sec:design})
show it would already break single-expression uses like
\texttt{aliceAge.toString() + aliceAge.toString()} only under
sub-expression granularity, but its statement-level version survives the
corpus surprisingly well — the reason to reject it is principle (it
havocs locals a call provably cannot reach), not corpus breakage.
(2) An attempt to justify clearing optionality in W-Field failed on the
index-covered-write case and was withdrawn (rule keeps optionality).
(3) An attempt to state Lemma~\ref{lem:rewiden-sound} without the
aliased-locals havoc failed — the two-locals-one-object counterexample is
noted in its proof; the lemma is honest only with the ``apply W-Field to
every local carrying a narrowed $f$'' clause.
(4) Effect tracking (option ii) was evaluated and found to collapse to
havoc at every ambient/indirect call; recorded so it is not re-proposed
without an annotation surface.

\end{document}
